% appendices.tex - the appendices of v_cvm.tex (informative)
\infappendix{The speech codecs as a channel: measurements}{measurements}
\section{Method}
The measurements were made with the D-Modem project's tools
[b-D-Modem], in October 2026, as follows.
\begin{enumerate}[label=\alph*)]
\item \textbf{Codecs.}  AMR: the OpenCORE AMR-NB library (encoder and
  decoder); AMR-WB: the VisualOn AMR-WB encoder and the OpenCORE AMR-WB
  decoder.  Both are the implementations of the 3GPP reference code that
  Android carries; the AMR-WB pair was checked against the 3GPP reference
  code of \tgpp{26.173} on a comb signal (clause~\ref{sec:crosscheck}).
  EVS: the 3GPP reference code of \tgpp{26.442} (fixed point, version
  16.4.0), encoder and decoder, through files in the G.192 format.  The
  codecs ran at the modes of Table~\ref{tab:codecmodes}, without DTX
  unless stated, one encoding and decoding per direction.
\item \textbf{Signal interface.}  16~bit linear PCM at 8000~Hz (AMR, and
  AMR-WB through a 2:1 resampler with a 3400~Hz low pass, as a narrowband
  signal reaches a wideband codec), at 16000~Hz (AMR-WB and EVS as they
  are) or at 32000~Hz (EVS super wideband).  The test signals were at
  $-12$~dBFS.
\item \textbf{Modulations of the V series.}  Two D-Modem modems
  (\code{half\_open8}) called each other through \code{netsim}, the
  project's network simulator, with the codec on its line
  (\code{NET\_CODEC}), at each modulation of Table~\ref{tab:modsurvive}
  with the modems' default settings (V.42 and V.44 where the modulation
  carries them), and sent 10~s of data each way; a call passed when both
  modems reached CONNECT and the data arrived intact.
\item \textbf{Waveform probes.}  Test signals were generated, passed
  through the codec (\code{codec\_probe}, or the EVS programs) and
  analysed against the input, aligned by correlation: steady tones (the
  signal-to-noise ratio of the output against the input, with the best
  gain), binary FSK at 300~bit/s (bit errors of a non-coherent detector),
  DQPSK and 8-DPSK on one carrier (symbol errors and the root mean square
  phase error of the differential detector), and the combs of
  clause~\ref{sec:modulation} with the phases changing at random every
  symbol (per tone: symbol errors, the rms phase error, and the rms
  change of the tone's amplitude from one symbol to the next).  Each
  measurement is over 150 symbols or 800 bits; a tone counts as
  \emph{good} when its symbol error ratio is below $2 \cdot 10^{-2}$, and
  the raw rate given is that of the good tones.
\item \textbf{Frame erasures.}  Frames were erased at random, independently,
  at the ratio given, and the decoder told so (its concealment ran).
\end{enumerate}
The scripts (\code{measure/probe.py}, \code{probe3.py}, \code{matrix.sh})
and the generator of the tables below (\code{tables.py}) are in the
D-Modem repository with this document, so that the measurements can be
repeated and extended.

\begin{table}[htbp]
\caption{The codec modes measured}
\label{tab:codecmodes}
\small
\begin{tabularx}{\textwidth}{lX}
\toprule
\rowcolor{ituhead}\thd{Codec} & \thd{Modes (kbit/s)} \\
\midrule
AMR (\tgpp{26.090}), 8000~Hz & 4.75, 7.4, 7.95, 10.2, 12.2 \\
AMR-WB (\rec{G.722.2}), 16000~Hz & 6.6, 12.65, 15.85, 23.85; also behind a 2:1 resampler at 8000~Hz \\
EVS (\tgpp{26.445}), primary modes & NB (8000~Hz): 7.2, 9.6, 13.2, 24.4; WB (16000~Hz): 9.6, 13.2, 16.4, 24.4, 32, 48, 64, 96, 128; SWB (32000~Hz): 13.2, 24.4, 32, 48, 64, 96, 128 \\
EVS AMR-WB IO & 23.85 (16000~Hz) \\
\bottomrule
\end{tabularx}
\end{table}

\section{The modulations of the V series through the codecs}
Table~\ref{tab:modsurvive} gives, for each modulation and codec, the rate
of the CONNECT the modems reached (bold where the data then arrived
intact in both directions; with an \textsuperscript{e} where the data had
errors, which means the modulation carried no error correction, or that
V.42's detection phase failed and the modems ran without it) and a dash
where the modems did not connect.
\begin{table}[htbp]
\caption{D-Modem's modulations through the codecs: CONNECT rate (bold: data intact; e: data in error; --: no connection)}
\label{tab:modsurvive}
\centering\footnotesize
\resizebox{\textwidth}{!}{\tabmodsurvive}
\end{table}
V.21 (and Bell~103) connect through every mode and carry their 300~bit/s
without error through AMR at 12.2 and AMR-WB at 23.85; at the lower modes
their bits are in error (Table~\ref{tab:fsk}).  V.22 connects through AMR
12.2 and its LAPM comes up, so that its data arrives intact: 1200~bit/s.
V.22bis connects through AMR 12.2 and AMR-WB 23.85 (its handshake is
V.22's) but its 16-point constellation does not survive, nor does V.42's
detection on it.  V.32 and above do not complete their training; V.23's
1200~bit/s channel does, its 75~bit/s channel not.

\begin{table}[htbp]
\caption{Bit error ratio of 300~bit/s FSK through the codecs (8000~Hz interface)}
\label{tab:fsk}
\centering\small
\tabfsk
\end{table}

\section{Steady tones}
Table~\ref{tab:tones} gives the signal-to-noise ratio of a steady tone
after the codec, which is high (25 to 37~dB through AMR 12.2) at most
frequencies and low (10 to 18~dB) at a few: those whose period is a whole
number of samples at 8000~Hz (800, 1000, 1333, 1600, 2000, 2667~Hz) and
at 1500~Hz.  A steady tone is the easiest signal for the long-term
predictor once it has locked, and the hardest for the linear predictor,
whose line spectrum it cannot place; the values are no guide to a
modulated signal, whose phases the predictor must follow, and the
single-carrier and comb measurements below are.
\begin{table}[htbp]
\caption{SNR (dB) of a steady tone at $-12$~dBFS after the codec, by frequency (Hz)}
\label{tab:tones}
\centering\footnotesize
\resizebox{\textwidth}{!}{\tabtones}
\end{table}

\section{One carrier}
Table~\ref{tab:carrier} gives the symbol error ratio and the rms phase
error of DQPSK on one carrier at 1700~Hz, at 300 and 600~baud (600 baud
is V.22's).  Through AMR 12.2 and AMR-WB 23.85 the 600~baud signal has a
symbol error ratio near $10^{-3}$ with 14$^\circ$ rms; through AMR 7.4
and below, and AMR-WB 12.65, it does not pass.
\begin{table}[htbp]
\caption{DQPSK on one carrier at 1700~Hz: symbol error ratio / rms phase error}
\label{tab:carrier}
\centering\small
\tabcarrier
\end{table}

\section{The combs}\label{sec:combmeas}
Table~\ref{tab:combnb} gives, at the 8000~Hz interface, the tone symbol
error ratio, the rms phase error and the raw rate of the good tones for
the narrowband comb (200~Hz spacing, 400 to 3400~Hz, 16 tones) with the
20~ms and the 22.5~ms symbols, with 8-DPSK, with the shorter 12.5~ms
symbol, and for the 400~Hz comb; Table~\ref{tab:combwb} the same at the
16000~Hz interface for the wideband combs (400 to 6800~Hz).  The columns
``+10~ms'' are the 20~ms symbol with its boundaries 10~ms into the codec
frame rather than at it.
\begin{table}[htbp]
\caption{Combs at the 8000~Hz interface: tone SER / rms phase error / raw rate of the good tones (bit/s)}
\label{tab:combnb}
\centering\scriptsize
\resizebox{\textwidth}{!}{\tabcombnb}
\end{table}
\begin{table}[htbp]
\caption{Combs at the 16000~Hz interface (AMR-WB): tone SER / rms phase error / raw rate of the good tones (bit/s)}
\label{tab:combwb}
\centering\scriptsize
\resizebox{\textwidth}{!}{\tabcombwb}
\end{table}
What the combs show, and clause~\ref{sec:modulation} builds on:
\begin{itemize}
\item the 16-tone comb with DQPSK and the 20~ms symbol passes AMR 12.2 at a
  tone symbol error ratio of $10^{-3}$ (12$^\circ$ rms) and AMR-WB 23.85
  better; it does not pass AMR 7.4 and below, nor AMR-WB 6.6;
\item the position of the symbol in the codec frame makes no difference;
\item 8-DPSK on the same comb has a symbol error ratio of $6 \cdot 10^{-2}$
  through AMR 12.2: too high for the code; on the 400~Hz comb (8 tones)
  it passes, at 1050~bit/s, less than DQPSK on 16 tones;
\item the shorter 12.5~ms symbol raises the error ratio by an order of
  magnitude: the codec needs the integration time;
\item at the 16000~Hz interface through AMR-WB 23.85, 28 of the 34 tones
  are good and the tones above 6400~Hz are not: AMR-WB codes 6.4 to
  7~kHz by bandwidth extension, without the waveform;
\item the amplitude of a tone changes from one symbol to the next by
  1.5 to 3~dB rms through the ACELP codecs: amplitude carries no bits
  there.
\end{itemize}

\section{Tandem codecs and frame erasures}
Table~\ref{tab:combnb} also has AMR 12.2 followed by AMR 12.2 (a
mobile-to-mobile call with a transcoding gateway, or two codecs in a
call's path), and AMR-WB 23.85 followed by AMR 12.2 (a wideband terminal
calling a narrowband one): the comb's error ratio roughly doubles, and
the loading drops a few tones.  With 1 and 3 per cent of the frames
erased at random the comb's error ratio rises by about the erasure ratio
itself (an erased frame is an erased symbol); the code of
Annex~\ref{annex:rs} is dimensioned for this.

\section{Cross-check of the codec implementations}\label{sec:crosscheck}
The AMR-WB encoder and decoder used (VisualOn, OpenCORE) were run on the
wideband comb and compared with the 3GPP reference code of \tgpp{26.173}
on the same input; the result is in the text of
Table~\ref{tab:crosscheck}.
\begin{table}[htbp]
\caption{The AMR-WB implementations against the 3GPP reference code}
\label{tab:crosscheck}
\centering\small
\begin{tabularx}{\textwidth}{>{\raggedright\arraybackslash}X>{\raggedright\arraybackslash}p{5.2cm}}
\toprule
\rowcolor{ituhead}\thd{Comparison} & \thd{Result} \\
\midrule
\tabcrosscheck
\bottomrule
\end{tabularx}
\end{table}

\infappendix{The EVS codec}{evs}
% evs.tex - Appendix: the EVS codec, measured with the 3GPP reference code
\section{The codec}
The enhanced voice services codec (EVS, \tgpp{26.445}) is the codec of
VoLTE where both terminals have it, and the codec of 5G voice (VoNR,
[b-GSMA NG.114]).  It codes 20~ms frames at 5.9 to 128~\kbs{} in four
bandwidths (narrowband to 4~kHz, wideband to 8~kHz, super wideband to
16~kHz, full band to 20~kHz), choosing per frame between an ACELP core
(speech), a transform core (MDCT: TCX and high-quality modes, for music
and mixed content) and, at the low rates, a generic sound coding of the
spectral envelope with a coarse residual; above the core's band it codes
a bandwidth extension that carries no waveform.  It also has an AMR-WB
interoperable mode (the AMR-WB codec, bit-exact) and a channel-aware mode
at 13.2~\kbs{} with partial redundancy.  Its signal classification
decides which core codes a frame, so what it does to a modem's signal
depends on what the signal looks like to it, not on the bit rate alone.

\section{Measurements}
The combs of clause~\ref{sec:modulation} were passed through the 3GPP
reference encoder and decoder of \tgpp{26.442} (fixed point, version
16.4.0, through G.192 files, DTX off) at the signal interface rates and
modes of Table~\ref{tab:codecmodes}, as Appendix~\ref{app:measurements}
describes.  Tables~\ref{tab:combevsnb} to \ref{tab:combevsswb} give, per
comb, the tone symbol error ratio, the rms phase error, the raw rate of
the good tones, and the rms change of a tone's amplitude from one symbol
to the next in dB (the last for the question of amplitude modulation,
clause~\ref{sec:qam}).

\begin{table}[htbp]
\caption{EVS at the 8000~Hz interface (narrowband): tone SER / rms phase error / raw rate of the good tones (bit/s) / rms amplitude change (dB)}
\label{tab:combevsnb}
\centering\scriptsize
\resizebox{\textwidth}{!}{\tabcombevsnb}
\end{table}
\begin{table}[htbp]
\caption{EVS at the 16000~Hz interface (wideband); EVS AMR-WB IO is the codec's AMR-WB mode}
\label{tab:combevswb}
\centering\scriptsize
\resizebox{\textwidth}{!}{\tabcombevswb}
\end{table}
\begin{table}[htbp]
\caption{EVS at the 32000~Hz interface (super wideband): the S200, S300 and S400 combs to 14~kHz}
\label{tab:combevsswb}
\centering\scriptsize
\resizebox{\textwidth}{!}{\tabcombevsswb}
\end{table}

\section{What they show}
\begin{enumerate}[label=\alph*)]
\item \textbf{At 7.2 and 9.6~\kbs{}} (narrowband) and \textbf{at 9.6 and
  13.2~\kbs{} wideband and super wideband}, the combs do not pass: the
  tone symbol error ratio is $10^{-1}$ or worse on every comb, and the
  amplitude of a tone changes by 5 to 9~dB rms from one symbol to the
  next.  The codec's low-rate generic coding keeps the spectral envelope
  and not the phases.  The robust mode (300~bit/s FSK) passes: V.21's
  bit error ratio is below $10^{-2}$ at every mode but 9.6 wideband.
  A single carrier at 600~baud has a symbol error ratio of $3 \cdot 10^{-2}$
  at 13.2 wideband: V.22 would connect and lose data.
\item \textbf{At 13.2~\kbs{} narrowband}, the N400 comb passes with a
  tone symbol error ratio below $10^{-3}$ (8$^\circ$ rms, 800~bit/s raw),
  the N200 comb loads 10 of its 16 tones, and 8-DPSK does not hold.  This
  is the mode a VoLTE call falls to in poor coverage; the probing's
  choice of the sparse comb is what keeps the DCE connected through it.
\item \textbf{At 24.4~\kbs{}}, the transform core takes over for the
  comb and the picture changes: narrowband, every comb passes with 8-DPSK
  (N200: $2 \cdot 10^{-3}$, 6$^\circ$, 2400~bit/s raw); wideband, W300
  passes DQPSK at $10^{-3}$ (2200~bit/s) and W400 8-DPSK on 11 of 17
  tones; super wideband, the band above about 8~kHz is the codec's
  bandwidth extension and carries nothing, and the probing loads 13 to
  15 tones (1300 to 1500~bit/s) -- less than the wideband interface gives
  at the same rate, because the signal's power is spread over tones the
  codec drops.  A DCE at a 32000~Hz interface should therefore offer the
  wideband combs as well (Table~\ref{tab:combs}: the probing finds the
  band).
\item \textbf{At 48~\kbs{} and above}, the codec carries the comb
  nearly as a waveform: wideband, 8-DPSK on all 33 tones of W200 at
  $3 \cdot 10^{-3}$ (4950~bit/s raw) and on all tones of W400 without
  error; super wideband at 48~\kbs{}, DQPSK on 31 of the 35 tones of S400
  to 12.4~kHz (3100~bit/s); at 96~\kbs{}, DQPSK on all 69 tones of S200
  (6900~bit/s) and 8-DPSK on 56 of them (8400~bit/s); at 128~\kbs{},
  8-DPSK on all 69 tones without error, 3.4$^\circ$ rms: 10\,350~bit/s raw.
\item \textbf{Amplitude.}  Through the ACELP codecs a tone's amplitude
  changes by 1.5 to 3~dB rms from one symbol to the next, and amplitude
  carries no bits.  Through EVS at 24.4~\kbs{} the change is 0.7 to
  1.0~dB rms, at 48 it is 0.5 to 1.0~dB, at 96 and 128 it is 0.3 to
  0.6~dB: two amplitude levels 6~dB apart, or the 16-point constellation
  of clause~\ref{sec:qam} with the trellis code, would hold at 48~\kbs{}
  and above.  With 16-QAM on the 69 tones of S200 the raw rate would be
  13\,800~bit/s, at 128~\kbs{} of codec rate; the loading is for further
  study, as that clause says.
\item \textbf{AMR-WB IO} at 23.85~\kbs{} behaves as the AMR-WB codec
  (which it is): W300 loads 20 of 23 tones (2000~bit/s), W200 25 of 33.
\item \textbf{Frame erasures} cost more through EVS than through AMR:
  with 1 per cent of the frames erased at 24.4~\kbs{} wideband, the
  comb's tone symbol error ratio rises from $10^{-3}$ to $2 \cdot 10^{-2}$
  (Table~\ref{tab:combevswb}), about twice the erasure ratio, because the
  transform core's overlap carries each erasure into the next frame as
  well: an erased frame costs two symbols.  The receiver's erasure
  marking (clause~\ref{sec:erasure}) should mark both, and the probing
  should leave erased symbols out of its statistics, which are over a
  few tens of symbols.
\item \textbf{Mode changes.}  A VoLTE call switches between 13.2 and
  24.4~\kbs{} with the radio conditions, and a VoNR call may hold
  higher modes.  Between 13.2 and 24.4 the usable loading changes from
  one comb (N400, 800~bit/s) to another with 8-DPSK (N200, 2400~bit/s):
  the renegotiation of clause~\ref{sec:reneg} is what follows the
  codec, and the DCE should expect it every few seconds in a moving
  terminal.
\end{enumerate}

\section{Consequences for the signal set}
The signal set of clause~\ref{sec:modulation} needs no change for EVS:
the combs, the two modulation orders and the probing's loading cover it
from 13.2~\kbs{} narrowband (N400, DQPSK) to 128~\kbs{} super wideband
(S200, 8-DPSK), a range of 800 to 10\,350~bit/s raw.  What EVS adds is
the case for the super wideband combs (S200 to S400), which this
Recommendation defines, and for the amplitude loading of
clause~\ref{sec:qam}, which it leaves for further study with the
measured amplitude stability as its starting point.  The terminal's
part is in Appendix~\ref{app:terminal}: a terminal that can hold the EVS
mode at 48~\kbs{} or more for a DCE's call multiplies its throughput.


\infappendix{Terminals and networks}{terminal}
This appendix says what a terminal and a network should do so that a DCE
of this Recommendation gets the most of the speech path; none of it is a
requirement of this Recommendation on them.
\section{The terminal}
\begin{enumerate}[label=\alph*)]
\item \textbf{Signal interface.}  The DCE needs the codec's PCM
  interface: in a UE, the point where the acoustic front end's output
  enters the encoder and the decoder's output leaves for the acoustic
  back end.  The acoustic processing of \tgpp{26.131} (echo cancellation,
  noise suppression, automatic gain control, equalization) is for the
  acoustic path and should be bypassed for the DCE's samples; noise
  suppression in particular treats a comb as noise.
\item \textbf{Codec selection.}  Where the terminal negotiates the codec
  (an IMS client, through SDP), it should offer, for a call the DCE will
  use, the highest modes available and the wideband codecs first: EVS at
  24.4~\kbs{} or more, AMR-WB at 23.85, AMR at 12.2; it should restrict
  the mode set (the \code{mode-set} parameter of [b-IETF RFC 4867]) to
  the highest modes, and set the EVS parameters \code{br} and \code{bw} of
  \tgpp{26.445} Annex~A accordingly.
\item \textbf{DTX and VAD.}  The terminal should turn discontinuous
  transmission off for the call (the \code{dtx} parameter of EVS; for AMR
  and AMR-WB it is the encoder's setting).  A comb is not speech, and a
  voice activity detector may classify it as noise and send comfort noise
  in its place.  The DCE's signals are continuous while a call is up, so
  that DTX, where it cannot be turned off, has nothing to detect as
  silence.
\item \textbf{Jitter buffer.}  The jitter buffer management of
  \tgpp{26.114} adapts its delay by dropping or repeating frames, or by
  time-scaling the decoded signal (\tgpp{26.448} for EVS).  Each
  adaptation is a step in the DCE's symbol timing of up to one frame,
  which the receiver re-acquires (clause~\ref{sec:timing}); a terminal
  should prefer, for a DCE's call, a fixed delay large enough for the
  path's jitter.
\item \textbf{Level.}  The terminal should pass the DCE's samples to the
  encoder at their level, without gain, and give the DCE the decoder's
  output at its level.
\end{enumerate}
\section{The network}
\begin{enumerate}[label=\alph*)]
\item \textbf{Transcoding.}  Each codec in the path costs throughput
  (Appendix~\ref{app:measurements}); transcoder-free operation between
  terminals with the same codec is best.  A gateway to the PSTN
  transcodes to \rec{G.711}, which costs nothing further.
\item \textbf{Echo cancellers.}  Network echo cancellers of [b-ITU-T
  G.168] disable on the 2100~Hz tone with phase reversals that \rec{V.8}'s
  ANSam is; a DCE sends it as a matter of course.  Mobile-to-mobile IMS
  calls have no network echo canceller; the UE's acoustic echo canceller
  is bypassed with the rest of the acoustic processing.
\item \textbf{Codec mode control.}  The network changes the codec mode
  with the radio conditions; the DCE follows with renegotiations
  (clause~\ref{sec:reneg}).  A network that holds a call at the highest
  mode while a DCE is in it gives it the highest throughput; the DCE
  cannot ask for that, and this Recommendation does not define a way.
\item \textbf{In-band signalling.}  DTMF and other tones sent out of band
  ([b-IETF RFC 4733]) do not reach the speech path; the DCE's signals are
  not DTMF and are unaffected.  A network that detects the V.8 tones and
  switches a PSTN gateway to a voice-band data mode ([b-ITU-T V.152]) or
  to modem relay ([b-ITU-T V.150.1]) does so on its own leg; the mobile
  leg's codec remains, and remains the limit.
\end{enumerate}

\infappendix{Implementation notes}{impl}
A DCE of this Recommendation can be built from the parts of a V series
software modem, as the D-Modem project [b-D-Modem] has them:
\begin{itemize}
\item the V.8 procedures, with the information field of
  Annex~\ref{annex:v8} added to the CM and JM encoders and decoders;
\item the V.21 modulator and demodulator for the robust mode and the
  INFO and MP messages, and the HDLC framing of V.42 for their frames;
\item a multitone modulator (an inverse discrete Fourier transform of
  the tone phases, the cyclic prefix prepended) and demodulator (the
  correlation of each tone over the useful part, as the probe scripts of
  Appendix~\ref{app:measurements} do it), with the differential
  detection, the Gray mapping of Table~\ref{tab:gray} and the scrambler
  of V.34;
\item a Reed-Solomon encoder and decoder over GF($2^8$) for the code of
  Annex~\ref{annex:rs}, with erasure decoding; a block interleaver;
\item the LAPM, V.42bis and V.44 of the V.42 implementation, unchanged,
  with the detection phase disabled;
\item the V.250 command interpreter with the \code{+MS} value and the
  \code{+CVM} command of clause~\ref{sec:dte};
\item for Annex~\ref{annex:fax}, the T.30 session machine and the T.38
  packetizer of a T.38 gateway or an internet-aware fax implementation.
\end{itemize}
The computational load is small: the comb has at most 68 tones at
50~symbols/s, and the Reed-Solomon code at most 3300~bit/s.  The
probing's phase error estimates and the loading are the only parts not
found in a V series modem, and they are a few hundred lines.

In D-Modem, the natural place is a data pump (\code{dsp/}) behind the
existing pump interface, which the modem core (\code{core/}) already
drives through V.8 for V.34 and V.90, with the link layer (\code{link/})
above it, and \code{netsim}'s \code{NET\_CODEC} as the test line; the
interoperability test harness then runs the DCE against itself through
each codec mode as it runs the other modems today.
